\section{Related Work}
\label{sec:related_work}

\subsection{Parameter-Efficient Fine-Tuning (PEFT)}
To mitigate the extreme storage and computational cost of training and serving separate large-scale models for every downstream task, parameter-efficient fine-tuning (PEFT) has emerged as a dominant research area. Notable paradigms include Bottleneck Adapters~\cite{houlsby2019parameter}, Prefix Tuning~\cite{li2021prefix}, Prompt Tuning~\cite{lester2021power}, BitFit~\cite{zaken2021bitfit}, and Compacter~\cite{karimi2021compacter}, many of which are unified under collaborative ecosystems such as AdapterHub~\cite{pfeiffer2020adapterhub}. In particular, Low-Rank Adaptation (LoRA)~\cite{hu2021lora} has gained widespread adoption due to its lack of additional inference latency under single-task settings. Low-Rank decomposition updates can be folded into the base model weights at compile time. Subsequent extensions optimize dynamic budget allocation~\cite{zhang2023adalora}, dynamic search-free rank~\cite{valipour2023dylora}, and high-efficiency quantization~\cite{dettmers2023qlora}. Toward a unified understanding, He et al.~\cite{he2021towards} present a framework that mathematically connects adapters, prefix tuning, and LoRA as variants of the same functional formulation. However, when serving multiple tasks concurrently on a single resource-constrained device, storing and dynamically switching entire base-model weights for each task remains highly inefficient, motivating the search for multi-expert serving architectures.

\subsection{Static Weight-Space Model Merging}
Rather than maintaining separate models, static model-merging methods combine multiple task-specific neural networks fine-tuned from a shared pre-trained base into a single, multi-task model without further training. Early approaches used simple parameter averaging (e.g., Model Soups~\cite{wortsman2022model}) or Fisher-weighted merging~\cite{matena2021merging}. Recent advances, including Task Arithmetic~\cite{ilharco2022editing}, TIES-Merging~\cite{yadav2023ties}, DARE~\cite{yu2023language}, and Dare-to-Merge~\cite{yu2023dare}, address interference by pruning redundant parameter updates or resolving sign conflicts before merging, often unified under modern model-merging toolkits like MergeKit~\cite{gallego2024mergekit}. While highly successful in static multi-task setups, these approaches merge weights *before* deployment and use a single, frozen set of weights for all samples. As a result, they fail to adapt to dynamic, mixed-task deployment streams, leading to ``heterogeneity collapse'' because a single, static parameter compromise cannot satisfy conflicting activation requirements across distinct tasks.

\subsection{Dynamic Merging and MoE Routing}
Dynamic serving models address heterogeneous streams by selecting or blending pathways on-the-fly. This relates closely to Mixture-of-Experts (MoE) architectures~\cite{shazeer2017outrageously, fedus2022switch, riquelme2021scaling}, which employ parametric routing layers to route tokens or samples to task-specific sub-networks. However, classical MoEs require joint, high-resource training and suffer from severe routing instability during fine-tuning.
To dynamically serve modular, independently trained experts without expensive joint training, non-parametric dynamic merging methods have been developed. Specifically, Parameter-Free Subspace Routing (PFSR)~\cite{pfsr2025} routes features by computing cosine similarity against classification weights. To prevent heterogeneity collapse, Micro-Batch Homogenization (MBH)~\cite{mbh2025} dynamically partitions heterogeneous batches into homogeneous micro-batches.
Additionally, dynamic modular composition can be achieved directly in activation space. For instance, AdapterFusion~\cite{pfeiffer2020adapterfusion} uses an attention-based fusion layer to dynamically combine independent task-specific adapters, while VL-Adapter~\cite{sung2022vl-adapter} extends parameter-efficient transfer learning to vision-and-language tasks.
While mathematically elegant, MBH introduces a substantial systems bottleneck: running separate forward passes for each active expert sequentially. This $O(K)$ execution scaling is highly inefficient on sequential processors. 

We also distinguish our paradigm from systems-level dynamic PEFT serving frameworks like S-LoRA~\cite{sheng2023slora}, Punica~\cite{chen2023punica}, and LoraHub~\cite{huang2024lorahub}. These frameworks are engineered for high-throughput multi-LoRA serving on large-scale GPU clusters by paging, scheduling, and composing discrete adapter weights in memory. While highly effective at managing multi-tenant LLM traffic, they rely on heavy system scheduling overhead, custom CUDA paging kernels, and high-resource server-class environments. They present a classic systems trade-off: they minimize memory bandwidth bottlenecks by caching active adapters, but at the cost of high scheduling complexity and significant GPU compute and VRAM footprints. In contrast, SPS-ZCA operates under a fundamentally distinct paradigm: instead of dynamically switching, paging, or scheduling discrete adapters in memory, it merges task-specific activations on-the-fly inside the shared neural layers in a single pass. This is completely training-free, compiler-friendly, and lightweight enough for resource-constrained edge CPUs, dramatically reducing scheduling complexity and compute footprints at the edge. Thus, our proposed SPS-ZCA implements single-pass activation blending ($O(1)$) to achieve perfect task-specificity without batch partitioning, scheduling overhead, or sequential hardware penalties.

\subsection{On-Device Edge Intelligence and Systems-ML}
Deploying deep learning on resource-constrained platforms (e.g., mobile devices, wearables, and microcontrollers) requires balancing model accuracy with deployment latency, memory, and energy constraints~\cite{lane2015deep, zhou2019edge}. The field of TinyML~\cite{warden2019tinyml} emphasizes that theoretical FLOP reductions often do not translate to actual latency savings due to memory bandwidth bottlenecks and hardware dispatch overheads. Our work adopts this Systems-ML co-design philosophy. We demonstrate that by avoiding batch splitting and sequential backbone invocation, we can bypass scheduling overhead and maximize hardware utilization on sequential edge CPUs.
